\part{Storage and the filesystem}

\chapter{Talking to the disk}

\begin{goals}
  \item What LBA addressing is and why it replaced CHS
  \item How ATA PIO transfers a sector
  \item Why the status register must be polled, and in what order
\end{goals}

\section{LBA: one number instead of three}

Chapter 13 used CHS, because that is what the BIOS offers in real mode. Once in protected mode we
talk to the disk controller directly, and we can use \term{LBA} (Logical Block Addressing): sectors
are numbered $0, 1, 2, \ldots$ and that is the whole model. No geometry, no off-by-one on sector
numbers, no track boundaries.

28-bit LBA addresses up to $2^{28}$ sectors, which is 128 GB. More than enough.

\section{The registers}

The ATA controller is a set of I/O ports:

\begin{center}
\begin{tabularx}{\textwidth}{@{}l l X@{}}
\toprule
\textbf{Port} & \textbf{Name} & \textbf{Purpose} \\
\midrule
\hexa{1F0} & DATA & 16-bit data transfer window \\
\hexa{1F2} & SECTOR\_COUNT & How many sectors to transfer \\
\hexa{1F3} & LBA\_LOW & LBA bits 7--0 \\
\hexa{1F4} & LBA\_MID & LBA bits 15--8 \\
\hexa{1F5} & LBA\_HIGH & LBA bits 23--16 \\
\hexa{1F6} & DRIVE & Drive select + LBA bits 27--24 \\
\hexa{1F7} & STATUS / COMMAND & Read: status. Write: command \\
\hexa{3F6} & ALT\_STATUS & Status without side effects \\
\bottomrule
\end{tabularx}
\end{center}

Note that \hexa{1F7} is two different registers depending on whether you read or write it. That kind
of thing is normal in hardware and catches people out constantly.

\section{Reading a sector}

\filetag{lytlnybl/kernel/drivers/ata.c}
\begin{lstlisting}[style=c]
bool ata_read_sector(uint32_t lba, void* buffer) {
  uint16_t *data = (uint16_t*)buffer;

  if (lba > 0x0FFFFFFF) return false;      /* beyond 28-bit LBA */

  ata_select_drive(lba);                   /* drive + top 4 LBA bits */

  outb(ATA_SECTOR_COUNT, 1);               /* one sector */

  outb(ATA_LBA_LOW,  lba & 0xFF);
  outb(ATA_LBA_MID,  (lba >> 8)  & 0xFF);
  outb(ATA_LBA_HIGH, (lba >> 16) & 0xFF);

  outb(ATA_COMMAND, ATA_CMD_READ_PIO);     /* 0x20 */

  if (!ata_wait_drq()) return false;       /* wait for data to be ready */

  for (int i = 0; i < 256; i++) {
    data[i] = inw(ATA_DATA);               /* 256 words = 512 bytes */
  }

  ata_400ns_delay();
  return true;
}
\end{lstlisting}

The 28-bit LBA is split across four registers because each port is 8 bits wide: three full bytes
plus the low nibble of the drive register.

\begin{concept}{Why 256 iterations and not 512}
The data port is 16 bits wide, so each \cod{inw} brings in two bytes. $256 \times 2 = 512$: one
sector. This is why the buffer is cast to \cod{uint16\_t*} --- and also why passing an odd-aligned
buffer would be a bad idea.
\end{concept}

\subsection{Waiting correctly}

\begin{lstlisting}[style=c]
static bool ata_wait_drq() {
  uint8_t status;
  for (;;) {
    status = inb(ATA_STATUS);
    if (status & ATA_STATUS_BSY) continue;   /* still busy: keep waiting */
    if (status & ATA_STATUS_ERR) return false;
    if (status & ATA_STATUS_DF)  return false;   /* device fault */
    if (status & ATA_STATUS_DRQ) return true;    /* data ready */
  }
}
\end{lstlisting}

\begin{warn}{The order of the checks is not optional}
BSY must be tested \emph{first}. While BSY is set, every other bit in the status register is
meaningless --- the controller is mid-update. Checking ERR before BSY can read a stale error bit
from a previous command and abort a perfectly good read.

Note also that this loop has no timeout. A disk that never clears BSY hangs the kernel forever.
Real drivers count iterations and give up.
\end{warn}

\begin{lstlisting}[style=c]
static void ata_400ns_delay() {
  inb(ATA_ALT_STATUS);
  inb(ATA_ALT_STATUS);
  inb(ATA_ALT_STATUS);
  inb(ATA_ALT_STATUS);
}
\end{lstlisting}

The specification requires 400 nanoseconds between certain operations. Each read of the alternate
status port takes about 100 ns on the bus, so four reads is the canonical way to wait. Reading the
\emph{alternate} status port matters: reading the normal one at \hexa{1F7} clears the pending
interrupt, which you may not want.

\section{Writing}

Writing mirrors reading with three differences: command \hexa{30} instead of \hexa{20},
\cod{outw} instead of \cod{inw}, and a cache flush at the end:

\begin{lstlisting}[style=c]
  outb(ATA_COMMAND, ATA_CMD_WRITE_PIO);
  if (!ata_wait_drq()) return false;
  for (int i = 0; i < 256; i++) outw(ATA_DATA, data[i]);
  ata_400ns_delay();
  outb(ATA_COMMAND, ATA_CMD_FLUSH_CACHE);   /* 0xE7 */
  ata_wait_bsy();
\end{lstlisting}

\begin{concept}{Why the flush}
Disks buffer writes internally and report success immediately. If power is lost before the buffer
reaches the platter, the data is gone. \cod{FLUSH\_CACHE} forces it out and does not return until it
is physically committed.

This is exactly what \cod{fsync()} does on a real system, and why databases care so much about it.
\end{concept}

% ============================================================
\begin{recipe}{reading and writing raw sectors}{ch37-ata-sectors}
Writes 512 bytes to LBA 2000 --- chosen to be far past both the kernel and the filesystem --- reads
them back into a different buffer and compares. Then it reads disk sector 70 and finds
\hexa{DEADBABE} in the first four bytes: the filesystem's superblock, seen as raw bytes by a driver
that knows nothing about filesystems.

Everything in the next three chapters is built on these two functions.
\end{recipe}

% ============================================================
\chapter{Designing a filesystem}

\begin{goals}
  \item What structures are needed to turn a disk into files
  \item What an inode is and why the name is not in it
  \item How the on-disk layout is organised
\end{goals}

\section{The problem}

A disk is an array of numbered 512-byte sectors. To turn that into ``files with names, in
directories, of arbitrary size'', you need to answer three questions and store the answers
\emph{on the disk itself}:

\begin{enumerate}
  \item Which blocks are free?
  \item For a given file, which blocks hold its contents, and how big is it?
  \item Given a name, which file is it?
\end{enumerate}

Each question gets its own structure: a \term{bitmap}, an \term{inode}, and a \term{directory}.

\section{The layout}

\begin{center}
\begin{tikzpicture}[font=\sffamily\scriptsize,x=1cm,y=1cm]
  \fill[redlight,draw=red]   (0,0)   rectangle (1.6,1.2);
  \node[align=center,font=\tiny] at (0.8,0.6) {super-\\block};
  \node[font=\ttfamily\tiny,text=gray] at (0.8,-0.3) {block 0};

  \fill[amberlight,draw=amber] (1.6,0) rectangle (3.2,1.2);
  \node[align=center,font=\tiny] at (2.4,0.6) {block\\bitmap};
  \node[font=\ttfamily\tiny,text=gray] at (2.4,-0.3) {block 1};

  \fill[accentlight,draw=accent] (3.2,0) rectangle (6.4,1.2);
  \node[align=center,font=\tiny] at (4.8,0.6) {inode table\\128 inodes};
  \node[font=\ttfamily\tiny,text=gray] at (4.8,-0.3) {blocks 2--13};

  \fill[greenlight,draw=green] (6.4,0) rectangle (12.0,1.2);
  \node[align=center,font=\tiny] at (9.2,0.6) {data blocks\\{\tiny file contents and directory contents}};
  \node[font=\ttfamily\tiny,text=gray] at (7.0,-0.3) {block 14 onwards};

  \draw[decorate,decoration={brace,amplitude=4pt},gray] (0,1.35) -- (12.0,1.35)
    node[midway,above,font=\sffamily\tiny,text=gray] {everything offset by FS\_START\_BLOCK = disk sector 70};
\end{tikzpicture}
\end{center}

\filetag{lytlnybl/include/filesystem/fs\_layout.h}
\begin{lstlisting}[style=c]
#define FS_BLOCK_SIZE 512
#define FS_START_BLOCK 70        /* the filesystem starts at disk sector 70 */
#define FS_SUPERBLOCK 0
#define FS_BITMAP_BLOCK 1
#define FS_INODE_START 2
#define FS_TOTAL_BLOCKS 1000
#define FS_TOTAL_INODES 128
#define FS_INODE_MAX_BLOCKS 10
#define FS_FILENAME_LENGTH 32
#define FS_MAGIC 0xDEADBABE
\end{lstlisting}

The whole filesystem is shifted 70 sectors into the disk, leaving room for the boot sector and the
kernel. The translation is one line:

\begin{lstlisting}[style=c]
bool fs_read_block(uint32_t block_num, void* buffer) {
  uint32_t sector = FS_START_BLOCK + block_num;
  return ata_read_sector(sector, buffer);
}
\end{lstlisting}

\section{The superblock}

Block 0 describes the filesystem itself:

\begin{lstlisting}[style=c]
typedef struct {
  uint32_t magic;          /* 0xDEADBABE */
  uint32_t block_size;
  uint32_t total_blocks;
  uint32_t bitmap_start;
  uint32_t inode_start;
  uint32_t inode_count;
  uint32_t inode_blocks;
  uint32_t root_inode;
  uint32_t data_start;
  uint32_t free_blocks;
  uint32_t free_inodes;
} fs_superblock_t;
\end{lstlisting}

\begin{concept}{The magic number}
\cod{fs\_read\_superblock} refuses to proceed if the magic does not match:
\begin{lstlisting}[style=c]
if (superblock->magic != FS_MAGIC) return false;
\end{lstlisting}
Without it, pointing the driver at an unformatted disk would make it read random bytes as
\cod{inode\_count} and \cod{data\_start}, and then happily write structures all over it. The magic
number costs four bytes and turns silent destruction into a clean failure.
\end{concept}

\section{Inodes}

\begin{lstlisting}[style=c]
typedef struct {
  uint32_t size;                        /* size in bytes */
  uint16_t type;                        /* free, file or directory */
  uint32_t blocks[FS_INODE_MAX_BLOCKS]; /* 10 direct block pointers */
} fs_inode_t;
\end{lstlisting}

\begin{concept}{Why the name is not in the inode}
This is the design decision that defines Unix-style filesystems, and it surprises everyone the first
time.

An inode describes a \emph{file}: its size, its type, where its data lives. The \emph{name} lives in
the directory that contains it. So a file can have several names (hard links), all pointing at one
inode and one copy of the data. Renaming a file touches only the directory, not the file.

lytlnyblOS does not implement hard links, but it inherits the structure --- which means adding them
later would be almost free.
\end{concept}

With 10 direct pointers and 512-byte blocks, the maximum file size is $10 \times 512 = 5120$ bytes.

\begin{warn}{5 KB is small, and the code knows it}
That is enough for the shell binary (4,392 bytes) but not much more --- and \cod{user\_test} plus
the shell together already use most of the inode table's useful capacity.

The classic solution is \term{indirect blocks}: the eleventh pointer points not at data but at a
block full of \emph{more} pointers. With 512-byte blocks that is 128 extra pointers, taking the
maximum to about 69 KB. A double-indirect pointer gets you to 8 MB. This is exactly how ext2 works,
and adding single indirection here is a very instructive exercise.
\end{warn}

\section{Directories are files}

A directory is an inode of type \cod{FS\_TYPE\_DIRECTORY} whose data blocks contain an array of
fixed-size records:

\begin{lstlisting}[style=c]
typedef struct {
  uint32_t inode;                  /* 0 means the slot is empty */
  char name[FS_FILENAME_LENGTH];   /* 32 bytes */
} fs_directory_entry_t;
\end{lstlisting}

Each entry is 36 bytes, so a 512-byte block holds 14 of them.

\begin{center}
\begin{tikzpicture}[font=\sffamily\scriptsize,
  b/.style={draw=grayborder,rounded corners=2pt,minimum width=3.0cm,minimum height=0.7cm,align=center}]
  \node[b,fill=accentlight] (i0) at (0,1.6) {inode 0\\{\tiny type=DIR, blocks[0]=14}};
  \node[b,fill=greenlight,minimum width=4.4cm] (d) at (5.4,1.6)
    {block 14 (directory contents)};
  \node[font=\ttfamily\tiny,align=left,anchor=west] at (3.4,0.55)
    {\{1, "bin"\}\\ \{2, "file.txt"\}\\ \{3, "notes.txt"\}\\ \{0, ""\} \ldots};
  \draw[-{Stealth[length=2mm]},accent] (i0) -- (d);
\end{tikzpicture}
\end{center}

\begin{concept}{Recursion bottoms out at the root}
To find \cod{/bin/shell} you read the root directory to find \cod{bin}, then read \cod{bin} to find
\cod{shell}. But how do you find the root, when there is no directory containing it?

By convention: the root is always inode 0. \cod{FS\_ROOT\_INODE} is the fixed point that makes the
whole recursive structure resolvable.
\end{concept}

\section{Path resolution}

\filetag{lytlnybl/filesystem/fs\_manager.c}
\begin{lstlisting}[style=c]
int32_t fs_resolve_path(const char* path) {
  if (!path || path[0] != '/') return -1;    /* absolute paths only */

  uint32_t current_inode = FS_ROOT_INODE;
  uint32_t i = 1;

  while (path[i] != '\0') {
    char name[FS_FILENAME_LENGTH];
    uint32_t name_length = 0;

    while (path[i] == '/') i++;              /* skip repeated slashes */
    if (path[i] == '\0') break;

    while (path[i] != '/' && path[i] != '\0') {
      if (name_length >= FS_FILENAME_LENGTH - 1) return -1;
      name[name_length++] = path[i++];
    }
    name[name_length] = '\0';

    int32_t inode = fs_find_directory_entry(current_inode, name);
    if (inode < 0) return -1;

    current_inode = inode;
  }
  return current_inode;
}
\end{lstlisting}

Split on slashes, look each component up in the current directory, move down. The bounds check
inside the inner loop is what stops a 200-character path component from overflowing the 32-byte
local array --- exactly the kind of check that is missing elsewhere in this codebase.

\begin{warn}{No \cod{.} or \cod{..}}
There is no current directory and no parent directory. Every path must be absolute, and
\cod{cd} does not exist. Adding them means creating two entries in every directory at \cod{mkdir}
time and handling them in resolution --- a contained and worthwhile exercise.
\end{warn}

% ============================================================
\chapter{File descriptors and the manager}

\begin{goals}
  \item What a file descriptor is and why it is just an index
  \item Why descriptors 0, 1 and 2 are special
  \item How reading and writing cross block boundaries
\end{goals}

\section{The open file table}

\filetag{lytlnybl/include/filesystem/fs\_manager.h}
\begin{lstlisting}[style=c]
#define FS_MAX_OPEN_FILES 16
#define FS_FD_OFFSET 3

typedef struct {
    uint32_t inode;
    uint32_t offset;    /* current read/write position */
    bool used;
} fs_file_t;
\end{lstlisting}

\begin{lstlisting}[style=c]
static fs_file_t open_files[FS_MAX_OPEN_FILES];

int32_t fs_open(const char* path) {
  int32_t inode = fs_resolve_path(path);
  if (inode < 0) return -1;

  for (int i = 0; i < FS_MAX_OPEN_FILES; i++) {
    if (!open_files[i].used) {
      open_files[i].used = true;
      open_files[i].inode = inode;
      open_files[i].offset = 0;
      return i + FS_FD_OFFSET;
    }
  }
  return -1;
}
\end{lstlisting}

\begin{concept}{A descriptor is an index, nothing more}
A file descriptor is not a pointer and not a handle. It is the array index of a slot in the kernel's
open-file table, plus 3.

That indirection is deliberate and important: the user program holds a small integer that means
nothing outside the kernel. It cannot forge a pointer to a file structure, cannot corrupt the table,
cannot reach another process's files. The kernel validates the index on every use:
\begin{lstlisting}[style=c]
fd -= FS_FD_OFFSET;
if (fd < 0 || fd >= FS_MAX_OPEN_FILES) return -1;
if (!open_files[fd].used) return -1;
\end{lstlisting}
This is the same design idea as virtual addresses: hand out names, not places.
\end{concept}

\subsection{Why the offset of 3}

Unix reserves three descriptors by convention:

\begin{center}
\begin{tabularx}{0.9\textwidth}{@{}l l X@{}}
\toprule
\textbf{fd} & \textbf{Name} & \textbf{In lytlnyblOS} \\
\midrule
0 & stdin  & The keyboard \\
1 & stdout & The screen \\
2 & stderr & The screen \\
\bottomrule
\end{tabularx}
\end{center}

So real files start at 3. That is what \cod{FS\_FD\_OFFSET} encodes, and it is why the system call
handler can dispatch on the number alone:

\filetag{lytlnybl/kernel/interrupts.c}
\begin{lstlisting}[style=c]
if (fd > 2) {
    regs->eax = fs_write(fd, (void*)buffer, count);   /* a real file */
    break;
}
/* otherwise: the screen */
vga_text_write(&terminal, (char*)buffer);
\end{lstlisting}

\begin{warn}{The table is global, not per process}
\cod{open\_files} is a single static array shared by every process. Two processes opening files get
different slots, which works --- but a process can pass any descriptor number and reach another
process's open file. And there is no cleanup: if a process exits without calling \cod{close}, the
slot leaks forever. Sixteen leaked descriptors and no file can be opened again.

A real system keeps a per-process descriptor table pointing into a global open-file table, and closes
everything on exit. That is two structures instead of one, and it is what makes descriptors actually
safe.
\end{warn}

\section{Reading across blocks}

A file's data lives in up to 10 scattered blocks, but the caller wants a flat byte range. The
translation is the interesting part:

\filetag{lytlnybl/filesystem/fs.c}
\begin{lstlisting}[style=c]
int32_t fs_read_file(uint32_t file_inode_num, void* read_buffer,
                     uint32_t size, uint32_t offset) {
  fs_inode_t file_inode;
  uint8_t buffer[FS_BLOCK_SIZE];

  if (!fs_read_inode(file_inode_num, &file_inode)) return -1;

  if (offset >= file_inode.size) return 0;      /* EOF is 0 bytes, not an error */

  if (offset + size > file_inode.size)
    size = file_inode.size - offset;            /* clamp to the real size */

  uint32_t bytes_read = 0;
  uint8_t* destination = (uint8_t*)read_buffer;

  while (bytes_read < size) {
    uint32_t position = offset + bytes_read;

    uint32_t block_index  = position / FS_BLOCK_SIZE;   /* which block */
    uint32_t block_offset = position % FS_BLOCK_SIZE;   /* where inside it */

    if (block_index >= FS_INODE_MAX_BLOCKS) break;

    if (!fs_read_block(file_inode.blocks[block_index], buffer)) return -1;

    uint32_t bytes = FS_BLOCK_SIZE - block_offset;      /* rest of this block */
    if (bytes > size - bytes_read)
      bytes = size - bytes_read;                        /* or less, if that is all we need */

    memcpy(destination + bytes_read, buffer + block_offset, bytes);
    bytes_read += bytes;
  }

  return bytes_read;
}
\end{lstlisting}

\begin{concept}{The same division as paging}
\cod{position / BLOCK\_SIZE} and \cod{position \% BLOCK\_SIZE}. It is exactly the arithmetic from
Chapter 26, where a virtual address split into a page number and an offset within the page.

The idea is identical: a flat linear space is mapped onto fixed-size chunks that live somewhere
scattered, and a division plus a remainder does the translation. Once you see this pattern, you see
it everywhere in systems --- pages, disk blocks, network packets, cache lines.
\end{concept}

\begin{warn}{EOF is 0, not an error}
The comment in the source says it: \cod{//EOF should return 0 bytes instead of an error}. This
matters enormously for the shell's \cod{cat}:
\begin{lstlisting}[style=c]
while ((n = read(fd, buffer, sizeof(buffer))) > 0) {
    write(1, buffer, n);
}
\end{lstlisting}
The loop ends on 0. If EOF returned $-1$, the condition \cod{> 0} would also end the loop --- but
the caller could not distinguish ``finished'' from ``disk error''. Getting these conventions right is
what makes a filesystem usable.
\end{warn}

% ============================================================
\begin{recipe}{creating, writing and reading files}{ch39-filesystem}
Makes a directory, creates a file in it, writes to it, closes and reopens it to read the contents
back, shows that a second read returns 0 for end of file rather than an error, lists both
directories and finally deletes the file.

The close-and-reopen is the point. A descriptor carries an offset that only moves forwards and
there is no \cod{lseek}, so after writing eight bytes the offset is at eight and reading returns
nothing. The kernel's own boot-time test does the same dance for the same reason.
\end{recipe}

% ============================================================
\chapter{mkfs: formatting from outside}

\begin{goals}
  \item Why a tool that runs on the host is needed
  \item How files get into the image before the system ever boots
  \item How the same header serves two completely different programs
\end{goals}

\section{The chicken and egg}

The system needs \cod{/bin/shell} to exist before it can run anything. But to create a file you need
a running system. Something has to break the cycle from outside.

\cod{mkfs} is that something: a program compiled for \emph{your} Linux machine, not for lytlnyblOS,
which opens the disk image as an ordinary file and writes the filesystem structures into it
directly.

\filetag{lytlnybl/Makefile}
\begin{lstlisting}[style=mk]
HOST_CC = gcc

$(MKFS): tools/mkfs.c
	mkdir -p $(dir $@)
	$(HOST_CC) -g -Iinclude $< -o $@
\end{lstlisting}

Note \cod{HOST\_CC = gcc}, not \cod{i686-elf-gcc}. This one program is deliberately built for the
host.

\begin{concept}{One header, two worlds}
\cod{mkfs.c} includes \filepath{include/filesystem/fs\_layout.h} --- the \emph{same} header the
kernel uses. So the definitions of the superblock, inodes and directory entries are shared between a
64-bit Linux program and a 32-bit freestanding kernel.

This is why the structures are so carefully defined with fixed-width types (\cod{uint32\_t}, not
\cod{int}). A plain \cod{int} might be a different size on each side, and the two programs would
disagree about the layout of the disk --- a bug that would look like random filesystem corruption.
\end{concept}

\section{What it does}

\begin{lstlisting}[style=c]
int main(int argc, char **argv) {
  FILE *image = fopen(argv[1], "r+b");     /* open the image read/write */

  printf("formatting %s..........\n", argv[1]);

  if (!format_filesystem(image)) { /* ... */ }

  fs_superblock_t sb_copy;
  uint8_t buffer[FS_BLOCK_SIZE];
  read_block(image, FS_SUPERBLOCK, buffer);
  memcpy(&sb_copy, buffer, sizeof(fs_superblock_t));

  if (!install_rootfs(image, &sb_copy)) { /* ... */ }

  fclose(image);
  printf("mkfs success\n");
}
\end{lstlisting}

Two phases: write empty structures, then copy in the contents of \filepath{rootfs/}. The
host-side block functions are the exact analogue of the kernel's:

\begin{lstlisting}[style=c]
static bool write_block(FILE *image, uint32_t block_num, const void *buffer) {
  uint32_t sector = FS_START_BLOCK + block_num;
  if (fseek(image, sector * FS_BLOCK_SIZE, SEEK_SET) != 0) return false;
  return fwrite(buffer, 1, FS_BLOCK_SIZE, image) == FS_BLOCK_SIZE;
}
\end{lstlisting}

\cod{fseek} plus \cod{fwrite} here; \cod{outb} plus \cod{inw} in the kernel. Same logical operation,
two completely different mechanisms, one shared layout.

\section{The build pipeline}

\filetag{lytlnybl/Makefile}
\begin{lstlisting}[style=mk]
$(KERNEL_IMG): $(BOOTLOADER) $(KERNEL_BIN) $(MKFS) $(USER_TEST_ROOTFS) $(SHELL_ROOTFS)
	dd if=/dev/zero of=$@ bs=512 count=20480   # a blank 10 MB image
	dd if=$(BOOTLOADER) of=$@ conv=notrunc     # boot sector at sector 0
	dd if=$(KERNEL_BIN) of=$@ seek=1 conv=notrunc  # kernel from sector 1
	$(MKFS) $@                                 # filesystem from sector 70
\end{lstlisting}

\begin{center}
\begin{tikzpicture}[font=\sffamily\scriptsize,
  n/.style={draw=grayborder,rounded corners=2pt,minimum width=2.5cm,minimum height=0.75cm,align=center}]
  \node[n,fill=purplelight] (a) at (0,0) {bootloader.asm};
  \node[n,fill=accentlight] (b) at (0,-1.1) {kernel/*.c, *.asm};
  \node[n,fill=greenlight] (c) at (0,-2.2) {user/*.c + libc.a};
  \node[n,fill=graylight] (d) at (0,-3.3) {rootfs/bin/*};

  \node[n,fill=amberlight,minimum height=3.6cm,minimum width=2.2cm] (img) at (5.6,-1.65) {\textbf{kernel.img}\\{\tiny 10 MB}};

  \draw[-{Stealth[length=2mm]},purple] (a) -- node[above,font=\tiny]{dd} (img);
  \draw[-{Stealth[length=2mm]},accent] (b) -- node[above,font=\tiny]{dd seek=1} (img);
  \draw[-{Stealth[length=2mm]},green] (c) -- node[below,font=\tiny]{cp to rootfs} (d);
  \draw[-{Stealth[length=2mm]},gray] (d) -- node[below,font=\tiny]{mkfs} (img);

  \node[n,fill=redlight,minimum width=2.6cm] (q) at (10.2,-1.65) {QEMU};
  \draw[-{Stealth[length=2.5mm]},red,thick] (img) -- node[above,font=\tiny]{make run} (q);
\end{tikzpicture}
\end{center}

\begin{tryit}
Run \cod{make} and watch the last lines:
\begin{lstlisting}[style=txt]
formatting build/kernel.img..........
mkfs success
\end{lstlisting}
Then boot. The kernel prints the contents of \cod{/} and \cod{/bin} during its self-test, and you
should see \cod{shell} and \cod{user\_test} listed --- files that were placed there by a program
running on your Linux machine, read back by a filesystem driver running on a 32-bit kernel you
wrote.
\end{tryit}

\begin{summary}
Storage works end to end: a disk driver, a filesystem, descriptors, and a way to get files in from
outside. Everything the shell needs now exists. The next part assembles it.
\end{summary}
